\chapter{Кооперативное отслеживание изолинии неизвестного скалярного поля группой неголономных морских роботов}\label{ch:ch1}

В данной главе рассматривается задача отслеживания неизвестной изолинии скалярного поля группой морских мобильных роботов. Каждый робот управляется через регулирование предельной угловой скорости вращения валов гребных винтов, причем выбор значений этих параметров ограничен физическими интервалами: скоростные режимы гребных винтов имеют минимальные и максимальные пороги. Динамика движительной системы определяется постоянными времени электромеханических приводов, коэффициентами тяги винтов, а также гидродинамическими эффектами, включая присоединенную массу, демпфирование, восстанавливающие силы и сопротивление поперечного потока, которые вычисляются с использованием эмпирических и теоретических коэффициентов. Система управления полностью автономна. Роботы не взаимодействуют друг с другом, каждый из них самостоятельно генерирует управляющие сигналы на основе собственных сенсорных данных.
В частности, каждый робот учитывает текущее значение поля в своей позиции, а также относительное положение и ориентацию соседних роботов в зоне видимости.
Важно подчеркнуть, что количество роботов в группе неизвестно ни одному из них, а их анонимность не позволяет идентифицировать отдельных участников.

Цель алгоритма — вывести группу на целевую изолинию с заданным значением поля, обеспечить устойчивое отслеживание и предотвратить кластеризацию, стремясь к квазиравномерному распределению. Сложность распределения вдоль неизвестной кривой обусловлена невозможностью прямого вычисления расстояний между роботами вдоль изолинии, а также двумерным характером движения, что исключает применение классических методов самораспределения агентов вдоль одномерных структур.

На защиту выносятся результаты, связанные с групповыми аспектами  алгоритма управления: отсутствие столкновений, предотвращение  кластеризации, компенсация влияния течений и эффективное самораспределение роботов вдоль изолинии.
Ключевым результатом главы является доказательство отсутствия столкновений и кластеризации за счет оригинального управления линейной скоростью, подавляющего боковое схождение роботов к общей траектории.
Для случая замкнутой изолинии доказана асимптотическая равномерность распределения. Теоретические выводы, включая анализ влияния гидродинамических сил и ограничений движительной системы, подтверждены результатами численного моделирования с учетом нелинейностей в динамике гребных винтов и пространственного взаимодействия роботов. Особое внимание уделено устранению дисбаланса между скоростными режимами винтов и внешними возмущениями, что критично для сохранения устойчивости при маневрировании в условиях переменного поля.

\section{Отслеживание изолинии отдельным роботом}\label{sec:ch1/sec1}

Для морского робота катамаранного типа, с постоянной скоростью движения $v$, вектором курса $\vec{e}(\psi) = [\cos(\psi), \sin(\psi)]^T$ и угловой скоростью $\omega$ измененная динамика с учетом потока воды $\vec{V}$ может быть описана следующим образом:
\begin{equation}
    \begin{cases}
        \dot{\vec{r}} = v\vec{e}(\psi) + \vec{V}(\vec{r}, t), \\
        \dot{\psi} = \omega,  \omega \in [\underline{\omega}; \overline{\omega}]
    \end{cases}
\end{equation}

где ${\dot{\vec{r}}}$ представляет результирующий вектор скорости, а $v$ - скорость относительно среды в моменте. $\vec{V}(\vec{r}, t)$ - скорость частицы среды в точке $r$ в момент времени $t$, $\psi$ - угол ориентации. $\vec{v}(\psi) = v\vec{e}(\psi)$ - скорость с заданными минимальными и максимальными значениями $\overline{v} > \underline{v} \geq 0$ Это обеспечивает адаптацию траектории робота к воздействию потока и инерции. Движение робота  начинается в положении $x(0) = x_0$, $y(0) = y_0$, $\psi(0) = \psi_0$.

Пусть есть некоторое скалярное поле интенсивности $F(\vec{r}) \in C^1(\mathbb{R}^3)$  и некоторое целевое значение $f_\ast$ (см. рисунок \ref{fig:scheme_circle}, где значение поля пропорционально интенсивности  серого цвета). Цель управления выражается $F(\vec{r}(t)) \underset{t \to \infty}{\longrightarrow} f_\ast$. Иными словами, существует некоторая целевая изолиния $I(f_\ast) = \{ \vec{r}: F(\vec{r})=f_\ast\}$, к которой катамаран должен ассимптотически стремиться. Системе управления доступно значение поля $f(t) := F(x,y)$ в текущей  локации робота $x = x(t)$, $y = y(t)$ и она способна определять скорость $f'(t)$, с которой меняются эти значения с течением времени $t$, например, используя  численное дифференцирование показаний датчика значения поля.

\begin{figure}[ht]
    %\begin{minipage}[b]{0.49\linewidth}\centering
    %    \includegraphics[width=0.5\linewidth]{scheme_circle.png}
    %\end{minipage}
    \centering
    \begin{minipage}[b]{0.49\linewidth}\centering
        \includegraphics[width=\linewidth]{scheme_circle.png}
    \end{minipage}
    \caption{Схематичное изображение поля искомой линии и движения ОУ}
    \label{fig:scheme_circle}
\end{figure}
$\vec{V}(\vec{r}, t)$ — скорость частицы среды $r$ в момент времени $t$, имеющая класс $C^1$;
$\vec{A}(\vec{r}, t)$ — ускорение частицы среды $r$ в момент времени $t$, является производной Лагранжа от скорости:
\begin{equation}
        \vec{A}(\vec{r}, t) = \vec{V}'_r + \vec{V}'_r \vec{V}
\end{equation}

$\Omega(\vec{r}, t)$ — завихренность среды:
\begin{equation}
    \frac{1}{2}[\vec{V}'_r - (\vec{V}'_r)^T] = \begin{bmatrix}
                                                    0 & -\Omega \\
                                                    \Omega & 0
                                                \end{bmatrix}
\end{equation}

$E(\vec{r}, t)$ — тензор скоростей деформации среды:
\begin{equation}
    E = \frac{1}{2}[\vec{V}'_r - (\vec{V}'_r)^T]
\end{equation}
$\bigtriangledown F(t, r)$ — пространственный градиент поля;
$I_{t,r} = I_t(F(t, r))$ — динамическая пространственная изолиния, проходящая через точку r в момент времени t;

\subsection{Постановка задачи и закон управления}\label{subsec:ch1/sec1/subsec1}
\subsection{Необходимые допущения и утверждения}\label{subsec:ch1/sec1/subsec2}

\subsection{Необходимые допущения и утверждения}\label{subsec:ch1/sec1/subsec3}

\section{Ссылки}\label{sec:ch1/sec2}

\section{Ссылки}\label{sec:ch1/sec3}

\section{Ссылки}\label{sec:ch1/sec4}

